% ===========================================================================
%  sviluppatore/06-dashboard/11-banco.tex
% ===========================================================================
\section{Studio › Blueprint \& code: il banco di lavoro}\label{sez:studio-banco}

\file{templates/studio.html}, \file{static/css/studio.css},
\file{static/js/studio.js}; la lavagna è \file{static/js/node\_editor.js}
con l'ossatura in \file{templates/\_v2\_studio\_bp.html}; il codice di
strategie e approcci è \file{templates/editor.html} in un iframe
(\file{static/js/editor.js}, \file{static/css/editor.css}).

La pagina è una superficie alta quanto la finestra, come Agent e Genetics
(\vediSez{studio-pagine}): \cod{.studio.wb}, una griglia con la colonna del
\emph{che cosa} a sinistra e il lavoro a destra. La classe \cod{wb} c'è solo
nella vista \cod{authoring}: la stessa pagina, servita come
\enquote{Dati esterni} alla pagina Data (\cod{data-view="data"}), resta un
modulo senza il banco.

\subsection{La colonna}

In alto le famiglie, \cod{.wb-type} con \cod{data-tab} uguale alla chiave
del pannello (\cod{strat}, \cod{approccio}, \cod{wm}, \cod{overlay},
\cod{metric}) e il conto nel sistema. Sotto, per ogni famiglia, una sezione
\cod{.inst-bar[data-inst]}: ricerca, \enquote{+}, importazione ed elenco.
Un solo modulo di \file{studio.js} le disegna da
\percorsoHTTP{/studio/api/installati/<categoria>} e le offre al resto della
pagina:

\begin{center}
\begin{tabular}{@{}L{40mm}L{98mm}@{}}
\toprule
\textbf{Funzione} & \textbf{Che cosa fa} \\
\midrule
\rigaTenue
\cod{\_\_instReload[cat]} & \cod{(nome, vista)}: ricarica l'elenco; con il
  nome lo segna aperto nella vista (la usano i salvataggi). \\
\cod{\_\_instApri[cat]} & \cod{(nome, vista)}: apre nella vista chiesta o in
  quella corrente; restituisce una promessa. \\
\rigaTenue
\cod{\_\_instSegna[cat]} & \cod{(nome, vista)}: segna l'aperto senza
  ricaricare (la usa l'editor nell'iframe). \\
\cod{\_\_studioNuovo[cat]} & \cod{(nome)}: weight mode, overlay e metriche
  ripartono dal template col nome dato. \\
\rigaTenue
\cod{studioActivateTab} & \cod{(chiave)}: apre una famiglia, cioè voce,
  elenco e pannello insieme. \\
\bottomrule
\end{tabular}
\end{center}

\noindent
\enquote{+} non crea niente finché non si preme Invio: il nome si scrive in
una riga in cima all'elenco e si controlla lì (\cod{[a-z][a-z0-9\_]*}, nomi
già presi). Poi dipende dalla vista: sulla lavagna si svuota la lavagna e
le si dà il nome (\simbolo{studioNuovaLavagna}, che chiede conferma se ci
sono blocchi non salvati); nel codice di strategie e approcci l'iframe crea
la cartella col template e la apre; negli altri si riparte dal template.
Eliminare passa per \simbolo{window.showConfirm} e per
\percorsoHTTP{/studio/api/installati/<categoria>/<nome>/elimina}.
La famiglia scelta e la colonna ridotta stanno in \cod{localStorage}
(\cod{qs-studio-tab}, \cod{qs-studio-rail}); la colonna ridotta la rimette
uno script in linea prima del primo disegno.

\subsection{Due viste, due documenti}

Strategie e approcci hanno la lavagna (\cod{.bp-pane}) e il codice
(\cod{.code-pane}), ognuna con la sua testa \cod{.wb-head} e la sua copia
dell'interruttore Blueprint/Code, nello stesso punto: \file{studio.js} tiene
allineate le due copie (\cod{aria-pressed}), passa il fuoco a quella della
vista che si apre e manda l'evento \cod{studio:modo}. A schermo intero va la
vista intera, testa compresa, quindi i comandi restano a portata.

Le due viste sono due documenti: possono tenere aperti due artefatti
diversi. La colonna ricorda l'aperto per vista (\cod{graph}, \cod{code}) e
segna quello della vista visibile. Un clic apre nella vista corrente; sulla
lavagna solo ciò che ha un grafo (\cod{blueprint} nell'elenco), il resto
passa al codice e la colonna lo dice.

\begin{perche}{Non si riapre da sé nell'altra vista}
Passare da una vista all'altra non carica l'artefatto aperto nell'altra:
sostituirebbe senza chiedere un grafo o un file con modifiche non salvate.
Il documento di ciascuna vista cambia solo per un gesto fatto lì.
\end{perche}

\subsection{La lavagna: Check e Save}

\simbolo{wireNodeCompile} collega i due comandi. \enquote{Check} chiama
\percorsoHTTP{/studio/api/compile} e mostra l'esito sotto la lavagna, con il
codice generato chiuso in un \cod{details}. \enquote{Save} compila e, se il
controllo passa e il codice ha il descrittore, installa
(\percorsoHTTP{/studio/api/install}, \percorsoHTTP{/studio/api/approach/install})
con il grafo accanto al codice; poi ricarica l'elenco con il nome. Ogni
\cod{onChange} della lavagna scrive \enquote{Unsaved changes};
\simbolo{studioLoadGraph} e \simbolo{studioNuovaLavagna} caricano con
\cod{\_\_carica} acceso, così aprire non conta come modifica.

I blocchi sono larghi 176 pixel (\cod{NODE\_W}), e le porte stanno a
$y + 19 + 20i$ come prima: le misure in \file{studio.css} le rispettano.
L'esempio della lavagna vuota ha le colonne a 214 pixel l'una dall'altra:
entra tutto a 1440 pixel con la libreria aperta, e fra due colonne resta
posto per le etichette \enquote{yes}/\enquote{no}.

\subsection{Il codice di strategie e approcci, nell'iframe}

\file{editor.html} con \cod{?embed=1} è l'editor a più file. In Studio
(\cod{body.in-studio}) la sua barra è nascosta: la testa del pannello mostra
cartella, file e stato e ha Save e Save \& test, e i due documenti si parlano
con \cod{postMessage}, accettato solo dalla stessa origine:

\begin{center}
\begin{tabular}{@{}L{36mm}L{102mm}@{}}
\toprule
\textbf{Messaggio} & \textbf{Da chi, e che cosa porta} \\
\midrule
\rigaTenue
\cod{editor-pronto} & Dall'iframe, finita l'apertura: i comandi mandati
  prima aspettavano in coda. \\
\cod{editor-stato} & Dall'iframe a ogni cambio: cartella, file, testo e tono
  dello stato, Save e Save \& test attivi, sola lettura. \\
\rigaTenue
\cod{editor-comando} & Da Studio: \cod{salva}, \cod{prova}, \cod{apri}
  (con il nome), \cod{crea} (con il nome). \\
\cod{ricarica-elenco} & Da Studio dopo un'importazione o un'eliminazione:
  se l'artefatto aperto non c'è più, l'editor si svuota e lo dice. \\
\rigaTenue
\cod{artefatti-cambiati} & Dall'iframe dopo un salvataggio o una creazione:
  la colonna si ricarica. \\
\bottomrule
\end{tabular}
\end{center}

\subsection{Stato, esito, tastiera}

Lo stato in testa prende un tono, \cod{data-tone} (\cod{ok}, \cod{err},
\cod{busy}), che il CSS colora. Lo deduce dal testo un
\cod{MutationObserver}, così vale per tutti i punti della pagina che lo
scrivono; l'iframe lo dice da sé (\cod{data-tone-fisso}). L'esito dei
controlli (\cod{.result}) si mostra quando c'è qualcosa da dire e si
richiude con la X. Ctrl/Cmd+S preme il pulsante con \cod{data-salva} della
vista visibile, o il Save dell'iframe; dentro un dialogo non fa niente.

\paragraph{ui-system.css.} Dà a \cod{.sec}, \cod{.studio-panel},
\cod{.inst-bar}, \cod{.bp-main} e \cod{.vsc} fondo e raggio con
\cod{!important}: dentro \cod{.studio.wb} si tolgono con un solo
\cod{:is(\ldots)}. Ai campi dà fondo, bordo e raggio con selettori che
pesano sei classi: il nome nel percorso, la ricerca e i campi dei blocchi si
riscrivono con \cod{\#qs-canvas} davanti. I contenitori di pulsanti hanno
\cod{data-ui-group="none"}, e percorsi ed elenchi \cod{data-i18n-skip},
perché il dizionario non \enquote{traduca} nomi di file
(\cod{overlays} diventava \enquote{Overlays}).

\subsection{Le prove}

\simbolo{provaBanco} in \file{tests/e2e\_studio.js}, nei due temi: la pagina
che sta nella finestra, famiglie e conti, ricerca, apertura e tastiera, la
famiglia ricordata; \enquote{+} con nomi sbagliati e giusti, Ctrl+S,
l'esito, l'eliminazione; la lavagna (esempio che entra tutto, Check, codice
generato, Save, Clear, riapertura); il codice nell'iframe comandato dalla
testa (percorso, Save, creazione, artefatto eliminato mentre è aperto); i
collegamenti \cod{?open=}, \cod{?tab=metric\&output=chart},
\cod{?import=overlay}; la colonna ridotta; 820 pixel; la vista
\enquote{Dati esterni}. \cod{SOLO=banco} prova solo questa parte.
